\documentclass[11pt]{article}
\usepackage[a4paper,margin=24mm]{geometry}
\usepackage{amsmath,amssymb,mathtools}
\usepackage[hidelinks]{hyperref}
\newcommand{\C}{\mathbb C}
\newcommand{\R}{\mathbb R}
\title{The actual period receiver has no repeated complex root at sufficiently large finite period}
\author{Independent mathematical derivation}
\date{September 2026}
\begin{document}\maketitle

\paragraph{Original objects and provenance.}
This calculation applies the complete CJR receiver \cite{CJR} to
the actual moment family calculated in PIN1--14 \cite{PIN}.
The period matrix and all its coefficients are the original PCL1--10
\cite{PCL}; the finite-period domain is PZU1--4 \cite{PZ}.
The marked matrix and quartet remain FC5 and FC15a, with the
canonical ES reader and its human-source attributions retained in
the cumulative source \cite{FC}. The source and target norms below
are exactly ECC24--26 and EQR1--6, including the Gamma masses,
centres and both source orders \cite{FC,EQR}. No free moment
vector is substituted for a period.

Throughout, $0<\delta<1/2$, $\gamma>2$, and the actual nonzero
unit is on the phase $a=\bar a\ne0$. Its value is retained:
$U=aI$ cancels with its own inverse in the original conjugation.
Write $z=u^{-1}$ and retain the exact constants
\[
 \beta=2(\gamma^2-\delta^2),\quad
 \eta=(\delta^2+\gamma^2)^2,\quad
 B=\beta^2-9\eta=-5\gamma^4-26\delta^2\gamma^2-5\delta^4<0,
 \quad h=9/B,
 \quad C=\frac{125(4\delta\gamma)^4\beta^8B^4}{81^4}>0.
 \tag{PRD1}
\]
The letter $C$ here is PIN7's scalar, not CJR's order constant.
The exponential variable is $\xi$; it is distinct from $z$.

\section{The actual entire coefficients and their first two terms}
Here is a complete source of every coefficient used below. Let
$\mathcal R(z)=\sum_{n\ge0}R_nz^n$ have precisely the entries
\[
 (R_n)_{a r}=
 \sum_{\substack{p,l\ge0\\2p+4l=5n+r+1-a}}
 \frac{\beta^p\eta^l}{3^p p!l!}
 (-5)^{p+l-n}(a/5)_{p+l-n},
 \quad 1\le a\le4,\quad0\le r\le3.
 \tag{PRD2}
\]
All sums at a fixed $n$ are finite. PCL2--5 prove convergence
for every complex $z$, $\det\mathcal R=1$, and the entire
inverse $\operatorname{adj}\mathcal R$. In particular, with
$S_0=R_0^{-1}$,
\[
 S_n=-S_0\sum_{r=1}^nR_rS_{n-r},\qquad
 T_{j,n}=\sum_{r=0}^nS_rD_jR_{n-r},\quad
 D_j=\operatorname{diag}(\zeta^{-j},\zeta^{-2j},
                           \zeta^{-3j},\zeta^{-4j}),
 \tag{PRD3}
\]
where $\zeta=e^{2\pi i/5}$ and every $j=1,2,3,4$ is retained.
Let $Q$ be the complete matrix PIN3 and define, for $0\le k\le3$,
\[
 F_{j,k,n}=
 \sum_{a=0}^k\sum_{r=0}^n
 \frac{(T_{j,r}e_a)^{\mathsf T}Q(T_{j,n-r}e_{k-a})}
 {a!(k-a)!}.
 \tag{PRD4}
\]
These are exactly $[\xi^kz^n]e^{-\xi}f_j(\xi,z)/(4i\delta\gamma)$
by PIN3--5. Thus every coefficient of the four actual moment
functions used here is the convergent, fully specified sum
\[
\begin{split}
 \nu_m(z)&=\partial_\xi^m(e^{-4\xi}E_A(\xi,z))\big|_{\xi=0},\\
 [z^N]\nu_m(z)&=m!(4i\delta\gamma)^4
 \sum_{\substack{k_1+\cdots+k_4=m\\n_1+\cdots+n_4=N}}
          \prod_{j=1}^4F_{j,k_j,n_j},\qquad0\le m\le3,\\
 \mu_m(z)&=\sum_{r=0}^m\binom mr4^{m-r}\nu_r(z).
\end{split}\tag{PRD5}
\]
There are no independent amplitudes in these expressions.
The last line is exactly $E_A=e^{4\xi}\prod_j(4i\delta\gamma F_j)$.

PIN6--8, including their parity proof, give
\[
 \nu_0=Cz^8+O(z^{10}),\quad
 \nu_1=-2Chz^7+O(z^9),\quad
 \nu_2=4Ch^2z^6+O(z^8),\quad
 \nu_3=-6Ch^3z^5+O(z^7).
 \tag{PRD6}
\]
Every remainder is the convergent Taylor tail from PRD2--5.
Consequently the functions $M_m(z)=z^{m-8}\mu_m(z)$ have
removable singularities at zero and are entire. Their exact first
two Taylor coefficients are
\[
\begin{array}{c|rrrr}
 m&0&1&2&3\\\hline
 M_m(0)&C&-2Ch&4Ch^2&-6Ch^3\\
 M'_m(0)&0&4C&-16Ch&48Ch^2
\end{array}.
\tag{PRD7}
\]
For example $\mu_2=\nu_2+8\nu_1+16\nu_0$ and
$\mu_3=\nu_3+12\nu_2+48\nu_1+64\nu_0$ give the last
two columns. Parity in PRD6 supplies the absence of every other
term that could contribute to this table.

\section{The full cubic, its discriminant and a convergent-coefficient radius}
Keep $b(X)=\sum_{q=0}^3b_qX^q$ from the literal interpolation
CJR1--2; in particular
\[
 b_3=\frac{355\sqrt2\gamma-214\delta+
                   i(308\sqrt2\delta+210\gamma)}
                   {616\delta\gamma(\delta^2+\gamma^2)}\ne0.
 \tag{PRD8}
\]
Let $\kappa=a_{\rm amp}/2-4-1/2$ as in ECC22--27.
For precision the unchanged PZU radius on this phase is
\[
 d_j(z)=f_j(0,z)=\sum_{n\ge1}d_{j,2n}z^{2n},\quad
 d_{j,2}=4i\delta\gamma L_j\ne0,\quad
 C_j=\sum_{n\ge1}|d_{j,2+2n}|,\quad
 R_A=\max\{R_*,1,\max_j(2C_j/|d_{j,2}|)^{1/2}\}.
 \tag{PRD8a}
\]
Here $L_j$ is exactly PIN4 and $R_*$ is the original proved
five-orbit radius. The coefficients are PRD2--4 at $k=0$.
PZU1--4 proves their convergence and $v=0$ at every finite
$|u|\ge R_A$. The actual first receiver polynomial is
\[
 m_1(P,z)=\sum_{q=0}^3b_q\sum_{r=0}^q
          \binom qr\mu_{q-r}(z)(P+\kappa)^r.
 \tag{PRD9}
\]
For $z\ne0$ introduce the invertible coordinate map
$X=z(P+\kappa)$, with inverse $P=X/z-\kappa$, and define
\[
 H(X,z)=z^{-5}m_1(X/z-\kappa,z)=\sum_{r=0}^3A_r(z)X^r,
 \quad
 A_r(z)=\sum_{q=r}^3\binom qr b_qz^{3-q}M_{q-r}(z).
 \tag{PRD10}
\]
This identity keeps the entire original $m_1=z^5H$; it does
not replace its original scalar or coordinate. Each $A_r$ is
entire by PRD7. Their boundary polynomial and first derivative are
\[
\begin{split}
 H(X,0)&=Cb_3p(X),\qquad
 p(X)=X^3-6hX^2+12h^2X-6h^3=(X-2h)^3+2h^3,\\
 \partial_zH(X,0)&=C(b_2+12b_3)(X-2h)^2
       =\left(4+\frac{b_2}{3b_3}\right)\partial_XH(X,0).
\end{split}\tag{PRD11}
\]
Substitution of all eight entries of PRD7 into PRD10 proves
both complete coefficient identities.

For the actual polynomial $H$, define its full discriminant
by the finite expression
\[
 \Delta(z)=A_2^2A_1^2-4A_3A_1^3-4A_2^3A_0
                    -27A_3^2A_0^2+18A_3A_2A_1A_0
             =\sum_{n\ge0}\Delta_nz^n.
 \tag{PRD12}
\]
The polynomial $p$ has discriminant $-108h^6$, as follows
by translating it to $Y^3+2h^3$ and substituting its coefficients
in PRD12. Its leading coefficient is one. Therefore
\[
 \Delta_0=-108C^4b_3^4h^6\ne0,\qquad \Delta_1=0,
 \qquad
 \operatorname{Disc}_P m_1(P,z)=z^{26}\Delta(z).
 \tag{PRD13}
\]
For the middle equality, PRD11 says that the first-order
coefficient variation is exactly that of the translation
$Cb_3p(X+(4+b_2/(3b_3))z)$; a translation does not change
root differences or the leading coefficient. Thus its discriminant
has zero first derivative. Equivalently, substitution of PRD7
into the five monomials of PRD12 gives zero directly.
For the final equality, the discriminant of a cubic with leading
coefficient $a_3$ and roots $r_j$ is
$a_3^4\prod_{j<k}(r_j-r_k)^2$; this equals PRD12 by expansion.
Multiplication of a cubic by $z^5$ gives a factor $z^{20}$,
and substituting $X=z(P+\kappa)$ gives a factor $z^6$.
This proves the exponent $26$, including all original constants.

Here is a positive, explicitly specified radius using the complete
actual coefficient tails, not an unspecified sufficiently-small
parameter. Set
\[
 T_\Delta=\sum_{n\ge2}|\Delta_n|<\infty,\quad
 r_\Delta=\begin{cases}
 1,&T_\Delta=0,\\
 \min\{1,(|\Delta_0|/(2T_\Delta))^{1/2}\},&T_\Delta>0,
 \end{cases}
 \qquad r_D=\min\{R_A^{-1},r_\Delta\}>0.
 \tag{PRD14}
\]
Convergence follows because PRD10--12 are entire power series;
absolute convergence at a larger radius also proves convergence
of the displayed coefficient sum. PRD2--5, PRD10 and PRD12
specify every term by finite exact convolutions. For $|z|\le r_D$,
\[
 |\Delta(z)-\Delta_0|\le T_\Delta|z|^2
       \le|\Delta_0|/2,
 \qquad
 |\operatorname{Disc}_Pm_1|\ge54C^4|b_3|^4|h|^6|z|^{26}>0
 \quad(z\ne0).
 \tag{PRD15}
\]
Also the leading coefficient of $m_1$ is $b_3\mu_0(z)\ne0$
by PZU4. Thus for every actual finite period $|u|\ge r_D^{-1}$
on this unit phase the original receiver has exactly three simple
complex roots. In particular
\[
 (m_1(P,z),\partial_Pm_1(P,z))\ne(0,0)
              \quad\text{for every }P\in\C.
 \tag{PRD16}
\]
This evaluates and excludes every actual CJR23--25 candidate in
this period domain. It uses the actual family PRD5 throughout.

\section{All three moving roots as exact convergent functions}
Write $\varrho=2^{1/3}>0$, $\omega_\ell=e^{2\pi i\ell/3}$ and
\[
 X_\ell=h(2-\varrho\omega_\ell),\quad 0\le\ell\le2,
 \qquad d=\sqrt3\varrho|h|,
 \quad m=\min_\ell|X_\ell|>0,
 \quad \epsilon=\tfrac14\min\{d,m\},
 \quad R=\max_\ell|X_\ell|+\epsilon.
 \tag{PRD17}
\]
The three $X_\ell$ are exactly the distinct roots of $p$, every
pair has distance $d$, and none is zero because $\varrho<2$.
Let $a_{r,n}=[z^n]A_r(z)$ from PRD10 and put
\[
 T_r=\sum_{n\ge1}|a_{r,n}|<\infty,\quad
 T_H=\sum_{r=0}^3T_rR^r,\quad
 L_H=|Cb_3|\epsilon(d-\epsilon)^2>0,
 \quad
 r_H=\begin{cases}1,&T_H=0,\\
 \min\{1,L_H/(2T_H)\},&T_H>0,
 \end{cases}
 \quad r_{\rm root}=\min\{r_D,r_H,m/(4(1+|\kappa|)),1\}>0.
 \tag{PRD18}
\]
On $\Gamma_\ell=\{X:|X-X_\ell|=\epsilon\}$ the other two
root distances are at least $d-\epsilon$, hence
$|H(X,0)|\ge L_H$. For $|z|\le r_{\rm root}$,
$|H(X,z)-H(X,0)|\le T_H|z|\le L_H/2$ on every contour.
Thus no polynomial in the interpolation
$H(X,0)+t(H(X,z)-H(X,0))$, $0\le t\le1$, vanishes there.
Its number of zeros inside each circle, with multiplicity, is one.
For a direct proof of this polynomial form of Rouch\'e's count,
factor each polynomial over $\C$: the integral of its logarithmic
derivative divided by $2\pi i$ is the sum of the circle integrals
of $(X-r_j)^{-1}$, hence the integer number of roots inside.
The integral depends continuously on $t$ since the denominator
never vanishes on the circle. It is therefore constant and equals
one at $t=0$. This proves the claimed count and simplicity without
an assumed root-selection rule.

The unique root in each circle is the exact holomorphic function
\[
 \mathcal X_\ell(z)=\frac1{2\pi i}
       \int_{\Gamma_\ell}
             X\frac{\partial_XH(X,z)}{H(X,z)}\,dX,
 \qquad
 P_\ell(z)=\mathcal X_\ell(z)/z-\kappa
                    \quad(0<|z|<r_{\rm root}).
 \tag{PRD19}
\]
Factorization and the same circle integral show that the first
integral equals the single root, including its multiplicity one.
Uniform nonvanishing of the denominator permits integration of
the convergent local $z$ power series, proving holomorphy.
These are all three original receiver roots. Moreover
$|\mathcal X_\ell(z)|\ge3m/4$ and $|\kappa z|\le m/4$,
so $P_\ell(z)\ne0$, as required for the EQR collision path.

The full Taylor series is determined recursively, not only its
limiting point. If $\mathcal X_\ell=X_\ell+\sum_{n\ge1}x_{\ell,n}z^n$,
then
\[
 x_{\ell,n}=-\frac{[z^n]H(X_\ell+
                     \sum_{j=1}^{n-1}x_{\ell,j}z^j,z)}
                    {Cb_3p'(X_\ell)},\quad
 p'(X_\ell)=3(X_\ell-2h)^2\ne0,
 \quad x_{\ell,1}=-4-\frac{b_2}{3b_3}.
 \tag{PRD20}
\]
The coefficient of the newly inserted $x_{\ell,n}$ is exactly
$\partial_XH(X_\ell,0)$; the contour construction proves convergence
of this recursively specified series on $|z|<r_{\rm root}$.
In original coordinates this gives all three expansions
\[
 P_\ell(z)=\frac{h(2-2^{1/3}\omega_\ell)}z
       -\kappa-4-\frac{b_2}{3b_3}+O(z),\qquad
 \partial_Pm_1(P_\ell(z),z)
       =Cb_3p'(X_\ell)z^6+O(z^7).
 \tag{PRD21}
\]
The derivative equality follows from the exact PRD10 chain rule.
Every leading constant is nonzero and retains the original quartet.

\section{The four collisions in both complete original Gamma norms}
Fix any of these finite actual periods and a centre $P\ne0$.
PRD16 makes $U_1=(m_1(P,z),m'_1(P,z))^{\mathsf T}\ne0$.
The EQR1--6 calculation therefore applies to all four retained
choices $\epsilon_c\in\{0,1\}$ and both signs of
$\lambda^2=(2\epsilon_c-1)3P/22$. For the collision parameter
$t=w_0^2$, it gives in each unchanged original metric
\[
 w_0Y\longrightarrow U_1/\lambda,\qquad
 |t|^{1/2}\|Y\|_{G_{b,t}}
       \longrightarrow\frac{\sqrt{U_1^*G_{b,0}U_1}}{|\lambda|}>0,
                \quad b=U,W.
 \tag{PRD22}
\]
Thus the bounded-jet alternatives EQR7--10 are excluded for every
complex centre throughout the evaluated period domain.

At the three actual moving roots, there is a further complete
evaluation of both original metric constants. Write
\[
 b_\Gamma=s/2,\quad M_\Gamma=(2\pi)^{s/2},\quad
 k_j=(b_\Gamma)_j/j!,\quad
 L_0=a_{\rm amp}/2+b_2/(3b_3),\quad
 d_\ell=p'(X_\ell),
 \qquad s\in\{1,a_{\rm amp}\}.
 \tag{PRD23}
\]
Both source orders are retained, as are the centres
$a_{\rm amp}/2$ and $a_{\rm amp}/2-4$.
At $P=P_\ell(z)$, the exact constant in PRD22 is
\[
 \mathcal E_{b,\ell}(z)=
 \sqrt{\frac{22}{3|P_\ell(z)|}}\,
 |\partial_Pm_1(P_\ell(z),z)|\,
                  \sqrt{(G_{b,0})_{22}}.
 \tag{PRD24}
\]
Here index $22$ denotes the derivative coordinate of the two-by-two
Gram. This formula is exact for every retained sign and either
collision-root choice; $|\lambda|=\sqrt{3|P|/22}$.

We calculate its period behaviour from the full ECC24--26
covariance factors, not a substituted metric. At $v=0$ their
actual coefficients satisfy, directly from PRD6,
\[
 g_0=Cz^8+O(z^{10}),\quad
 g_1=2iChz^7+O(z^9),\quad
 g_2=-2Ch^2z^6+O(z^8),\quad
 g_3=-iCh^3z^5+O(z^7).
 \tag{PRD25}
\]
Let $e(P)=(1,P,P^2,P^3)$, let $Q_U,Q_W$ be the exact
covariance factors in ECC26, and put
$J_b=\left(\begin{smallmatrix}e(P)\\e'(P)\end{smallmatrix}\right)
Q_bQ_b^*
\left(\begin{smallmatrix}e(P)\\e'(P)\end{smallmatrix}\right)^*$,
$G_{b,0}=J_b^{-1}$.
For the source, multiplication by the complete upper triangular
$B_F$ in ECC25 gives the four entries of $e(P)Q_U$, after their
unchanged column factors $M_\Gamma^{-1/2}k_j^{-1/2}$, as
\[
\begin{split}
 c_0&=g_0,\qquad c_1=g_1-i(P-1/2)g_0,\\
 c_2&=g_2-b_\Gamma g_0/2-i(P-1/2)g_1
                                      -(P-1/2)^2g_0/2,\\
 c_3&=g_3-b_\Gamma g_1/2
       -i(P-1/2)(g_2-(3b_\Gamma+2)g_0/6)\\
    &\hspace{10mm}-(P-1/2)^2g_1/2+i(P-1/2)^3g_0/6.
\end{split}\tag{PRD26}
\]
The derivative row is the derivative of these same expressions
with $z$ fixed. Since
$z(P_\ell-1/2)=X_\ell-zL_0+O(z^2)$ by PRD21,
the entire two-row limit is
\[
 z^{-6}\begin{pmatrix}e(P_\ell)Q_U\\e'(P_\ell)Q_U\end{pmatrix}
 \longrightarrow
 \frac{Cd_\ell}{6\sqrt{M_\Gamma}}
 \begin{pmatrix}
 0&0&-k_2^{-1/2}&-iL_0k_3^{-1/2}\\
 0&0&0&i k_3^{-1/2}
 \end{pmatrix}.
 \tag{PRD27}
\]
To see the cancellation explicitly, the leading $z^5$ term in
$c_3$ is $iCp(X_\ell)z^5/6=0$; its next term is
$-iCd_\ell L_0z^6/6$. The leading $c_2$ is
$-C(X_\ell-2h)^2z^6/2=-Cd_\ell z^6/6$.
The derivative of $c_3$ has leading $iCd_\ell z^6/6$;
all other derivative entries have higher order. PRD25 proves
that the omitted actual Taylor tails enter only at those higher
orders, establishing every entry of PRD27.
The limiting two rows are independent. Taking their full Gram
and its inverse therefore gives
\[
 |z|^{12}(G_{U,0})_{22}\longrightarrow
 \frac{36M_\Gamma(k_3+k_2|L_0|^2)}{C^2|d_\ell|^2}.
 \tag{PRD28}
\]
For example the unscaled two-by-two Gram in PRD27 has entries
$k_2^{-1}+|L_0|^2k_3^{-1}$, $-L_0k_3^{-1}$,
$-\bar L_0k_3^{-1}$ and $k_3^{-1}$; its determinant is
$(k_2k_3)^{-1}$. This proves the displayed inverse entry and
retains every Gamma factor.

For the target $Q_W$ is independent of the period. Write $q_j$
for its row with index $j=0,1,2,3$. Polynomial expansion gives
$\|e(P)Q_W\|^2\sim|P|^6\|q_3\|^2$, while the determinant
of its two-row Gram is the squared exterior norm and satisfies
\[
 \det J_W\sim |P|^8\|q_3\wedge q_2\|^2,
 \quad
 |P|^2(G_{W,0})_{22}\longrightarrow
       \frac{\|q_3\|^2}{\|q_3\wedge q_2\|^2}
       =2M_\Gamma b_\Gamma(b_\Gamma+1).
 \tag{PRD29}
\]
Indeed the leading wedge is
$P^3q_3\wedge2Pq_2+P^2q_2\wedge3P^2q_3
=-P^4q_3\wedge q_2$.
The last two rows of the complete ECC26 factor are upper triangular;
their diagonal magnitudes are respectively
$[2M_\Gamma b_\Gamma(b_\Gamma+1)]^{-1/2}$ and
$[6M_\Gamma b_\Gamma(b_\Gamma+1)(b_\Gamma+2)]^{-1/2}$.
Their off-diagonal term affects neither the wedge nor the ratio,
proving the last equality rather than omitting that entry.

Combining PRD21, PRD24, PRD28 and PRD29 proves the two actual
positive period limits, for each $\ell$ and each original $s$:
\[
\begin{split}
 \lim_{z\to0}|z|^{-1/2}\mathcal E_{U,\ell}(z)
 &=6|b_3|\sqrt{\frac{22M_\Gamma(k_3+k_2|L_0|^2)}{3|X_\ell|}},\\
 \lim_{z\to0}|z|^{-15/2}\mathcal E_{W,\ell}(z)
 &=C|b_3||d_\ell|
     \sqrt{\frac{44M_\Gamma b_\Gamma(b_\Gamma+1)}
                         {3|X_\ell|^3}}.
\end{split}\tag{PRD30}
\]
These are iterated limits: first the EQR collision parameter tends
to zero at its fixed actual finite period; then the period tends
to infinity on the constructed branch. No simultaneous-limit
interchange or uniform collision assertion is assumed.
In particular the first marked receiver column can vanish at
each of its three moving simple roots, but its derivative there
remains nonzero and all four collision states escape in both
original quotient norms. Their full complements remain EQR3--4.
The result proves absence of repeated receiver roots on the
specified finite-period region. It does not assert an original
conductor order jump or an RH endpoint.

\begin{thebibliography}{9}
\bibitem{PIN} Split-Zero programme, \emph{The exact four singular
losses at the infinite-period boundary},
\texttt{PERIOD\_INVERSE\_BODY.tex}, PIN1--14.
\bibitem{CJR} \emph{The complete complex-centre jet map of the
original real-period receiver},
\texttt{COMPLEX\_RECEIVER\_JETS.tex}, CJR1--25 and CJR17a.
\bibitem{PCL} Split-Zero programme, \emph{The literal original
conductor limit and an invertible pivot restriction},
\texttt{PCL\_COMPLETE.tex}, PCL1--10 and PCL22--24.
\bibitem{PZ} Split-Zero programme,
\texttt{PERIOD\_ZERO\_COMPLETE\_PROOF.tex}, PZX2--3 and
PZU1--4. These earlier results already prove the actual phase's
order-two factor values and the finite domain $v=0$.
\bibitem{FC} Split-Zero programme,
\texttt{FABLE\_TO\_ORIGINAL\_CONDUCTOR.tex}, FC5, FC15a,
ECC22--27 and ECC39. Original human-source attribution to the
canonical ES reader, Levent Alp\"oge, Fable and Terence Tao is
retained there, as are the NIST DLMF sources for the original
Meixner--Pollaczek Gamma measure and norms.
\bibitem{EQR} Split-Zero programme,
\texttt{ESCAPE\_QUOTIENT\_RETURN\_BODY.tex}, EQR1--12.
\end{thebibliography}
\end{document}
